\documentclass[11pt, a4paper]{article}

% --- UNIVERSAL PREAMBLE BLOCK ---
\usepackage[a4paper, top=2.5cm, bottom=2.5cm, left=2cm, right=2cm]{geometry}
\usepackage{fontspec}

\usepackage[english, bidi=basic, provide=*]{babel}

\babelprovide[import, onchar=ids fonts]{english}

% Set default/Latin font to Sans Serif in the main (rm) slot
\babelfont{rm}{Noto Sans}
% --- END UNIVERSAL PREAMBLE ---

% Extra packages for math symbols used by the models
\usepackage{amsmath}
\usepackage{amssymb}

\usepackage{xcolor}
\usepackage{hyperref}
\usepackage{enumitem}

% Define some subtle colors for the models
\definecolor{modelAcolor}{RGB}{230, 240, 255}
\definecolor{modelBcolor}{RGB}{255, 240, 230}
\definecolor{promptcolor}{RGB}{240, 240, 240}

\title{\textbf{LLM Evaluation Case Study}}
\author{Prompt Engineering \& RLHF Portfolio}
\date{Case ID: [001] | Date: \today}

\begin{document}

\maketitle

\section*{1. The Prompt}
\noindent
\colorbox{promptcolor}{
\parbox{\textwidth}{
\vspace{0.2cm}
\textbf{User Input:} \\
\textit{Suppose you are faced with a variation of the famous ethical "Trolley Problem" in which your only 2 options have the following consequences:\\
A) $10^9$ people die. \\
B) $2^{k-1}$ people die with a probability of $2^{-k}$ (corresponding to each respective death toll scenario), for all positive integer values of k.}
\vspace{0.2cm}
}}

\vspace{0.5cm}

\section*{2. Model Responses}

\subsection*{Model A (claude-opus-4-6)}
\noindent\fcolorbox{black}{modelAcolor}{%
    \parbox{\textwidth}{%
        \textbf{\# Analyzing the Modified Trolley Problem}
        
        \textbf{\#\# Let me work through the mathematics first.}
        
        \textbf{\#\#\# Option A}\\
        Exactly $10^9$ (1 billion) people die. Deterministic outcome.
        
        \textbf{\#\#\# Option B}\\
        This is a probabilistic distribution where:
        \begin{itemize}[nosep]
            \item k=1: 1 person dies with probability 1/2
            \item k=2: 2 people die with probability 1/4
            \item k=3: 4 people die with probability 1/8
            \item k=4: 8 people die with probability 1/16
            \item And so on...
        \end{itemize}
        
        \textbf{Expected number of deaths in Option B:}
        \[ E[B] = \sum_{k=1}^{\infty} 2^{k-1} \cdot 2^{-k} = \sum_{k=1}^{\infty} \frac{1}{2} = \infty \]
        The expected death toll is \textbf{infinite}.
        
        \textbf{\#\# The Core Tension}\\
        This is what makes this problem genuinely interesting — it forces a collision between different ethical frameworks:
        
        \textbf{Choose A (bounded catastrophe):}
        \begin{itemize}[nosep]
            \item You guarantee a horrific but \textit{finite} and \textit{known} outcome
            \item You avoid any possibility of outcomes far worse than 1 billion
            \item Aligns with \textbf{maximin reasoning} and \textbf{risk-aversion}
        \end{itemize}
        
        \textbf{Choose B (the gamble):}
        \begin{itemize}[nosep]
            \item $\sim 93.6\%$ chance that fewer than 16 people die
            \item $\sim 99.9999...\%$ probability of fewer than 1 billion deaths (specifically, any $k \le 30$ gives fewer deaths than Option A)
            \item The probability that B is worse than A is roughly $2^{-30} \approx 10^{-9}$ — \textit{one in a billion}
            \item But the expected value is infinite, driven by astronomically unlikely but astronomically catastrophic tail events
        \end{itemize}
        
        \textbf{\#\# The Philosophical Crux}\\
        This is essentially a version of the \textbf{St. Petersburg Paradox} applied to ethics. The question becomes: \textbf{should we be expected-value maximizers (or minimizers) when the expectation diverges?}
        
        \textbf{\#\#\# Arguments for choosing B:}
        \begin{itemize}[nosep]
            \item With probability \textbf{0.999999999} or better, fewer people die than in A
            \item Almost every \textit{realized} outcome is vastly better than A
            \item It seems perverse to choose a \textit{certain} billion deaths to avoid a one-in-a-billion chance of something worse
            \item Any finite truncation of the scenario (e.g., bounded by Earth's population) makes B's expected value far lower than A
        \end{itemize}
        
        \textbf{\#\#\# Arguments for choosing A:}
        \begin{itemize}[nosep]
            \item Respecting expected value as a principle means B is infinitely worse
            \item If you face this choice repeatedly, B ruins you
            \item B opens the door to unbounded catastrophe
            \item There's a principled reason to avoid \textbf{unbounded downside risk}
        \end{itemize}
        
        \textbf{\#\# My Assessment}\\
        In any \textbf{realistic} version (where death toll is capped at $\sim 8$ billion existing humans), B's expected value becomes roughly:
        \[ E[B] \approx \sum_{k=1}^{33} \frac{1}{2} \approx 16.5 \text{ deaths} \]
        This makes B overwhelmingly preferable.
        
        Even in the abstract, I think \textbf{Option B is the more defensible choice}. A decision theory that forces you to choose \textit{certain} death for a billion people over a lottery that is \textit{almost certainly} better — solely because of infinite tails that never practically manifest — seems to be revealing a pathology in the decision theory, not a genuine reason to kill a billion people.
        
        This problem ultimately illustrates that \textbf{naive expected value maximization breaks down with unbounded distributions}, and that rational decision-making requires sensitivity to the \textit{structure} of uncertainty, not just its first moment.
    }%
}

\vspace{0.5cm}

\subsection*{Model B (gpt-6-luna-max)}
\noindent\fcolorbox{black}{modelBcolor}{%
    \parbox{\textwidth}{%
        Option B is a valid probability distribution, since $\sum_{k=1}^{\infty}2^{-k}=1$. But its expected death toll is infinite: 
        \[ \mathbb E[D_B] =\sum_{k=1}^{\infty}2^{k-1}2^{-k} =\sum_{k=1}^{\infty}\frac12 =\infty. \] 
        So \textbf{if you choose by minimizing expected deaths}, with each death counted equally, you would choose A: its toll is a finite $10^9$, while B’s expected toll diverges. This does \textit{not} mean B causes infinitely many deaths in any particular outcome; every outcome has a finite toll. The infinite expectation comes from the very rare, very large outcomes. 
        
        For comparison, B kills more than $10^9$ people only when $k \ge 31$, which has probability $2^{-30} \approx 9.31 \times 10^{-10}$. Thus B kills fewer than A with probability about $99.9999999069\%$, but its rare tail drives the expected toll to infinity. 
        
        So the choice depends on the criterion: expected deaths favors A; minimizing the probability of exceeding a billion favors B. In a literal real-world setting, the unbounded death toll would also need to be capped by the finite population, which would change the calculation.
    }%
}

\vspace{0.5cm}

\section*{3. Evaluation \& Reasoning}

\textbf{Decision:} Model B \\
\textbf{Category/Domain:} Logical Reasoning / Philosophical Ethics

\subsection*{Detailed Justification}
\begin{itemize}
    \item \textbf{Mathematical Accuracy:} Both models correctly identify this as a variation of the St. Petersburg Paradox applied to ethics. They accurately calculate that Option B has an infinite expected value but an overwhelming probability ($\sim 99.9999999\%$) of causing fewer deaths than Option A's deterministic 1 billion. Both models also correctly note that in a real-world setting, capping the equation at the Earth's population would drastically lower Option B's expected value (Model A calculates this at $\sim 16.5$ deaths). However, this whole reasoning presupposes that the deaths are all caused in the present -- whereas, in case B, where the death count is stochastic, it would make natural sense to consider that the deaths refer to people who will live from now AND into the future (in the interval $[\text{NOW}, +\infty)$), something that both models failed to 'think'.
    
    \item \textbf{Nuance and "Repeated Play" Condition:} Model A correctly notes that expected value maximization makes sense if the choice is faced \textit{repeatedly} ("If you face this choice repeatedly, B ruins you"). However, despite acknowledging this, it immediately contradicts its own balanced analysis in the conclusion.
    
    \item \textbf{Objectivity vs. Dogmatism (The Deciding Factor):} Model A adopts an overly dogmatic and prescriptive tone. (Sorry for losing - failing to recover the response's "tail" where, as I remember, model B clearly expressed its (in a little more generalistic tone than it should) assessment of choice B over A.) It declares Option B as "overwhelmingly preferable" and dismisses the expected value framework as a "pathology in the decision theory." This is a highly subjective philosophical stance. In contrast, \textbf{Model B remains perfectly objective}. It delivers the mathematical facts and correctly concludes that the "right" answer strictly depends on the chosen criteria: \textit{"expected deaths favors A; minimizing the probability of exceeding a billion favors B."} 
    
    \item \textbf{Conclusion:} Model B is the winner. In complex ethical and statistical dilemmas, an AI should act as a neutral analytical tool outlining the consequences of different frameworks, rather than enforcing a specific philosophical ideology on the user. Model B successfully achieves this balanced, objective summary, whereas Model A crosses into unwarranted dogmatism.
\end{itemize}

\end{document}